
%\documentclass[a4paper,12pt,landscape]{article}
\documentclass[a4paper,14pt,landscape]{extarticle}

\usepackage[a4paper,margin=8mm,bmargin=16mm]{geometry}
\usepackage[german,main=french]{babel}

\usepackage{fontspec}
\usepackage{microtype}

%\setmainfont{Libertinus Serif}
%\setmainfont{Junicode}
%\setmainfont{Vollkorn}
\setmainfont{Gentium}
% https://www.teuderun.de/schriftarten/top-10/

\usepackage{paracol}

% https://tex.stackexchange.com/a/505562/295527
\newcommand\chunks[2]{%
\begin{leftcolumn*}\begin{otherlanguage}{german}%
{#1}%
\end{otherlanguage}\end{leftcolumn*}%
\begin{rightcolumn}\begin{otherlanguage}{french}%
{#2}%
\end{otherlanguage}\end{rightcolumn}%
}

\usepackage{csquotes}
% \textelp{}
% https://tex.stackexchange.com/a/294634/295527

\begin{document}

\thispagestyle{empty}
\setlength{\columnsep}{8mm}
%\setlength{\columnseprule}{0.4pt}
\setlength{\columnseprule}{0pt}

\begin{paracol}{2}
  \chunks{\begin{center}
          \textbf{Die moderne Technik}
          \end{center}}
         {\begin{center}
          \textbf{La technique moderne}
          \end{center}}
         
  \chunks{Was ist die moderne Technik? Auch sie ist ein Entbergen. Erst wenn wir den Blick auf diesem Grundzugruhen lassen, zeigt sich uns das Neuartige der modernen Technik.}
         {Qu’est-ce que la technique moderne ? Elle aussi est un dévoilement. C’est seulement lorsque nous arrêtons notre regard sur ce trait fondamental que ce qu’il y a de nouveau dans la technique moderne se montre à nous.}
  \chunks{\textelp{}}
         {\textelp{}}
  \chunks{Das Wasserkraftwerk ist in den Rheinstrom gestellt. Es stellt ihn auf seinen Wasserdruck, der die Turbinen daraufhin stellt, sich zu drehen, welche Drehung diejenige Maschine umtreibt, deren Getriebe den elektrischen Strom herstellt, für den die Uberlandzentrale und ihr Stromnetz zur Strombeförderung bestellt sind. Im Bereich dieser ineinandergreifenden Folgen der Bestellung elektrischer Energie erscheint auch der Rheinstrom als etwas Bestelltes. Das Wasserkraftwerk ist nicht in den Rheinstrom gebaut wie die alte Holzbrücke, die seit Jahrhunderten Ufer mit Ufer verbindet. Vielmehr ist der Strom in das Kraftwerk verbaut. Er ist, was er jetzt als Strom ist, nämlich Wasserdrucklieferant, aus dem Wesen des Kraftwerks. Achten wir doch, um das Ungeheuere, das hier waltet, auch nur entfernt zu ermessen, für einen Augenblick auf den Gegensatz, der sich in den beiden Titeln ausspricht: \glqq Der Rhein\grqq, verbaut in das Kraftweih, und \glqq Der Rhein\grqq, gesagt aus dem Kunstwerk der gleichnamigen Hymne Hölderlins. Aber der Rhein bleibt doch, wird man entgegnen, Strom der Landschaft. Mag sein, aber wie? Nicht anders denn als bestellbares Objekt der Besichtigung durch eine Reisegesellschaft, die eine Urlaubsindustrie dorthin bestellt hat.}
         {La centrale électrique est mise en place dans le Rhin. Elle le somme de livrer sa pression hydraulique, qui somme à son tour les turbines de tourner. Ce mouvement fait tourner la machine dont le mécanisme produit le courant électrique, pour lequel la centrale régionale et son réseau sont commis aux fins de transmission. Dans le domaine de ces conséquences s’enchaînant l’une l’autre à partir de la mise en place de l’énergie électrique, le fleuve du Rhin apparaît, lui aussi, comme quelque chose de commis. La centrale n’est pas construite dans le courant du Rhin comme le vieux pont de bois qui depuis des siècles unit une rive à l’autre. C’est bien plutôt le fleuve qui est muré dans la centrale. Ce qu’il est aujourd’hui comme fleuve, à savoir fournisseur de pression hydraulique, il l’est de par l’essence de la centrale. Afin de voir et de mesurer, ne fût-ce que de loin, l’élément monstrueux qui domine ici, arrêtons-nous un instant sur l’opposition qui apparaît entre les deux intitulés : \frquote{Le Rhin}, muré dans l’usine d’énergie, et \frquote{Le Rhin}, titre de cette œuvre d’art qu’est un hymne de Hölderlin. Mais le Rhin, répondra-t-on, demeure de toute façon le fleuve du paysage. Soit, mais comment le demeure-t-il ? Pas autrement que comme un objet pour lequel on passe une commande, l’objet d’une visite organisée par une agence de voyages, laquelle a constitué là-bas une industrie des vacances.}

  \chunks{\begin{flushright}
          \large\textsc{Heidegger}, \textit{Die Frage nach der Technik}.
          \end{flushright}}
         {\begin{flushright}
          \large\textsc{Heidegger}, \textit{La question de la technique}.
          \end{flushright}}

\end{paracol}

\end{document}
